\documentclass[../../main.tex]{subfiles}
\begin{document}
% SOURCE p=201; printed=186; object=ch06-basic-exercises-title; region=lower
\subsection*{习题}
% SOURCE p=201; printed=186; object=ch06-ex01; number=1; region=lower
\begin{exercise}\label{exercise:ch06-minkowski-equality}
Minkowski 不等式何时取等号？（$p=1$ 与 $1<p<\infty$ 的答案不同；$p=\infty$ 时怎样？）
\end{exercise}
% SUPPLEMENT: id=v14-ch06-ex01-solution; parent=ch06-ex01; kind=solution+motivation+hypotheses; origin=independent-derivation
\begin{solution}
\textbf{全部指数与退化情形（重构新增）。}总假设 $f,g\in L^p$。
\begin{enumerate}
\item 当 $p=1$，等号恰当 $f(x)\overline{g(x)}$ 几乎处处为非负实数。事实上，非负可积函数 $|f|+|g|-|f+g|$ 积分为零当且仅当它几乎处处为零；而对复数 $z,w$，
\begin{align*}
(|z|+|w|)^2-|z+w|^2=2\bigl(|z||w|-\operatorname{Re}(z\overline w)\bigr),
\end{align*}
所以复数三角不等式的等号等价于 $z\overline w$ 非负实。这允许任一个值为零；不要求两个函数有固定的全局比例。
\item 当 $1<p<\infty$，若任一函数为零元素，则等号成立。若 $A=\|f\|_p>0$、$B=\|g\|_p>0$，等号恰当 $f/A=g/B$ 几乎处处，即二者为固定非负倍数。证明如下：令 $u=f/A$、$v=g/B$、$t=A/(A+B)\in(0,1)$，逐点有
\begin{align*}
|tu+(1-t)v|^p
\leqslant(t|u|+(1-t)|v|)^p
\leqslant t|u|^p+(1-t)|v|^p.
\end{align*}
实函数 $s\mapsto s^p$ 在 $[0,\infty)$ 严格凸，因为其导数 $ps^{p-1}$ 严格递增；故第二个不等式取等号当且仅当 $|u|=|v|$。再结合复数三角等号判据，两步同时等号当且仅当 $u=v$。Minkowski 等号使上述左、右两端积分都为1，非负差积分为零，所以逐点等号几乎处处成立。反之，若 $u=v$，则 $f+g=(A+B)u$，立刻取等号。包含零函数的统一表述是存在 $a,b\geqslant0$、$(a,b)\ne(0,0)$，使 $af=bg$ 几乎处处。
\item 当 $p=\infty$，记 $A=\|f\|_\infty$、$B=\|g\|_\infty$。若 $A=0$ 或 $B=0$，总取等号。若 $A,B>0$，准确判据为：对每个 $\varepsilon>0$，集合
\begin{align*}
E_\varepsilon=\{x:\;&|f(x)|>A-\varepsilon,\quad |g(x)|>B-\varepsilon,\\
&|f(x)|+|g(x)|-|f(x)+g(x)|<\varepsilon\}
\end{align*}
具有正测度。即两个模长须在同一批正测度点上同时接近其各自本质上确界，且相位抵消趋于零。若 $\|f+g\|_\infty=A+B$，则 $\{|f+g|>A+B-\varepsilon\}$ 有正测度；删去使 $|f|>A$ 或 $|g|>B$ 的零集，在余下集合上，
\begin{align*}
(A-|f|)+(B-|g|)+(|f|+|g|-|f+g|)<\varepsilon.
\end{align*}
三项非负，故此集合包含于 $E_\varepsilon$，证明必要性。反之，在 $E_\varepsilon$ 上有 $|f+g|>A+B-3\varepsilon$；正测度性给 $\|f+g\|_\infty\geqslant A+B-3\varepsilon$，令 $\varepsilon\downarrow0$ 并结合三角不等式即得等号。条件不受零集改值影响。在 $(0,1)$ 上取 $f=1,g=x$，就有等号，但二者不是固定倍数，且和不在任何点取到本质上确界2。
\end{enumerate}
\end{solution}
\begin{remark}
\textbf{动机与条件说明（重构新增）。}$p=1$ 只检验逐点无相位抵消；$1<p<\infty$ 的严格凸性再把相对大小固定下来；$p=\infty$ 只要求在接近峰值的正测度区域内几乎不损失。有限范数保证非负差可积或峰值有限；零范数须单独处理。端点判据不能要求最大值实际取得，也不能只找一个可能属于零集的点。
\end{remark}
% END SUPPLEMENT: id=v14-ch06-ex01-solution
% SOURCE p=201; printed=186; object=ch06-ex02; number=2; region=lower
\begin{exercise}\label{exercise:ch06-linfty-properties}
证明\cref{theorem:ch06-linfty-properties}。
\end{exercise}
% SUPPLEMENT: id=v14-ch06-ex02-solution; parent=ch06-ex02; kind=solution; origin=independent-derivation
\begin{solution}
\textbf{完整证明对应（重构新增）。}本题要求的五款论证已逐款写在\hyperref[proof:ch06-linfty-properties-complete]{该定理后的完整补证}中：第一款由本质上界及非负差积分给 Hölder 与等号；第二款证明范数全部公理；第三款构造统一满测集并证明反向；第四款构造可测有界的一致 Cauchy 极限；第五款用复平面有限网格给简单逼近。该处证明独立于本题的引用，覆盖全部五款与零值情形，故无循环。
\end{solution}
% END SUPPLEMENT: id=v14-ch06-ex02-solution
% SOURCE p=201; printed=186; object=ch06-ex03; number=3; region=lower
\begin{exercise}\label{exercise:ch06-lp-intersection-banach}
若 $1\leqslant p<r\leqslant\infty$，则 $L^p\cap L^r$ 在范数 $\|f\|=\|f\|_p+\|f\|_r$ 下是 Banach 空间；若 $p<q<r$，则包含映射 $L^p\cap L^r\to L^q$ 连续。
\end{exercise}
% SUPPLEMENT: id=v14-ch06-ex03-solution; parent=ch06-ex03; kind=solution+motivation+hypotheses; origin=independent-derivation
\begin{solution}
\textbf{范数、完备与包含连续性（重构新增）。}交空间对线性运算封闭；非负性、齐次性、三角不等式由两范数相加得到，而 $\|f\|=0$ 强制 $f=0$ 几乎处处。设 $(f_n)$ 在和范数中 Cauchy，则在两空间分别 Cauchy，由已证完备性，存在 $u\in L^p$、$v\in L^r$ 使 $\|f_n-u\|_p\to0$、$\|f_n-v\|_r\to0$。须验证两极限是同一等价类。

若 $r<\infty$，对任意 $\varepsilon>0$，由三角不等式的集合包含及积分估计，
\begin{align*}
\mu(|u-v|>\varepsilon)
&\leqslant\mu(|u-f_n|>\varepsilon/2)+\mu(|f_n-v|>\varepsilon/2)\\
&\leqslant(2/\varepsilon)^p\|u-f_n\|_p^p
 +(2/\varepsilon)^r\|f_n-v\|_r^r\longrightarrow0.
\end{align*}
若 $r=\infty$，充分大 $n$ 时 $\|f_n-v\|_\infty\leqslant\varepsilon/2$，第二集合为零集，第一项仍趋零。所以 $\mu(|u-v|>\varepsilon)=0$。取 $\varepsilon=1/k$ 的可数并，得 $u=v$ 几乎处处。记共同元素为 $f$，则 $f\in L^p\cap L^r$，并在和范数下 $f_n\to f$。

若 $p<q<r$，由插值估计，
\begin{align*}
\|f\|_q\leqslant\|f\|_p^\lambda\|f\|_r^{1-\lambda}
\leqslant\|f\|_p+\|f\|_r.
\end{align*}
最后一步在右边为零时直接成立，否则分别以右边控制两个因子。包含映射线性且有界，故连续。
\end{solution}
\begin{remark}
\textbf{动机与条件说明（重构新增）。}和范数要求两种控制同时收敛；关键不是分别得到两个极限，而是证明它们代表同一个可测函数。$p,r\geqslant1$ 使两空间为 Banach，$p<r$ 保证 $p$ 有限并可用水平集估计；$r=\infty$ 以本质界替代第二项积分。中间 $q$ 保证插值权重在 $(0,1)$，全证不依赖测度有限性。
\end{remark}
% END SUPPLEMENT: id=v14-ch06-ex03-solution
% SOURCE p=201; printed=186; object=ch06-ex04; number=4; region=bottom
\begin{exercise}\label{exercise:ch06-lp-sum-banach}
若 $1\leqslant p<r\leqslant\infty$，则 $L^p+L^r$ 在范数
\begin{align*}
\|f\|=\inf\{\|g\|_p+\|h\|_r:f=g+h\}
\end{align*}
下是 Banach 空间；若 $p<q<r$，则包含映射 $L^q\to L^p+L^r$ 连续。
\end{exercise}
% SUPPLEMENT: id=v14-ch06-ex04-solution; parent=ch06-ex04; kind=solution+motivation+hypotheses; origin=independent-derivation
\begin{solution}
\textbf{下确界范数的完整验证（重构新增）。}记题中量为 $N(f)$。分解集合非空，故 $N(f)$ 是有限非负数；改变几乎处处代表不改变分解关系及两范数。对 $c\ne0$，分解 $cf=cg+ch$ 与 $f=g+h$ 一一对应，给 $N(cf)=|c|N(f)$；$c=0$ 时同样成立。对任意 $\eta>0$，选 $f=g+h$、$f'=g'+h'$ 使费用分别小于 $N(f)+\eta$、$N(f')+\eta$；分解 $f+f'=(g+g')+(h+h')$ 并用两空间三角不等式，得到 $N(f+f')\leqslant N(f)+N(f')+2\eta$，令 $\eta\downarrow0$。

若 $N(f)=0$，可选 $f=g_n+h_n$ 且 $\|g_n\|_p+\|h_n\|_r<1/n$。对每个 $\varepsilon>0$，若 $r<\infty$，
\begin{align*}
\mu(|f|>\varepsilon)\leqslant(2/\varepsilon)^p\|g_n\|_p^p
 +(2/\varepsilon)^r\|h_n\|_r^r\longrightarrow0.
\end{align*}
若 $r=\infty$，充分大 $n$ 的 $|h_n|$ 几乎处处小于 $\varepsilon/2$，仅余第一项。因此 $f=0$ 几乎处处，证明正定性；这里没有假设下确界能取得。

为证完备，使用一般范数空间的绝对级数判据。设 $\sum_kN(f_k)<\infty$，取分解 $f_k=g_k+h_k$ 满足
\begin{align*}
\|g_k\|_p+\|h_k\|_r<N(f_k)+2^{-k}.
\end{align*}
于是两边分量分别组成 Banach 空间中的绝对可和级数，令 $g=\sum_kg_k\in L^p$、$h=\sum_kh_k\in L^r$。由定义，
\begin{align*}
N\left(\sum_{k=1}^nf_k-(g+h)\right)
\leqslant\left\|\sum_{k=1}^ng_k-g\right\|_p+
       \left\|\sum_{k=1}^nh_k-h\right\|_r\longrightarrow0.
\end{align*}
故 $N$ 范数完备。

最后，若 $f\in L^q$ 非零，令 $F=f/\|f\|_q$。按正文的模长1截断，$F=G+H$，其中 $\|G\|_p\leqslant1$、$\|H\|_r\leqslant1$（$r=\infty$ 时由处处界1）。故 $N(f)\leqslant2\|f\|_q$；零元素直接成立。包含映射因而连续。
\end{solution}
\begin{remark}
\textbf{动机与条件说明（重构新增）。}同一函数可能有许多分解，下确界选择总费用最小的极限值；无需存在最优分解。正定性须排除通过抵消而费用趋零的可能，水平集估计正承担此事。$p,r\geqslant1$ 提供分量的范数和完备性，$p<q<r$ 提供大值与小值分别受控的截断。可和误差 $2^{-k}$ 让逐项近似最优分解仍可求和，不要求测度空间有限。
\end{remark}
% END SUPPLEMENT: id=v14-ch06-ex04-solution
% SOURCE p=201,202; printed=186,187; object=ch06-ex05; number=5; region=bottom-top
\begin{exercise}\label{exercise:ch06-lp-inclusion-criterion}
设 $0<p<q<\infty$。则 $L^p\not\subseteq L^q$ 当且仅当 $X$ 含有任意小的正测度集合；$L^q\not\subseteq L^p$ 当且仅当 $X$ 含有任意大的有限测度集合。（对“若”方向：第一种情形存在不交序列 $\{E_n\}$，使 $0<\mu(E_n)<2^{-n}$；第二种情形存在不交序列 $\{E_n\}$，使 $1\leqslant\mu(E_n)<\infty$。对适当的常数 $a_n$ 考虑 $f=\sum_na_n\chi_{E_n}$。）$q=\infty$ 时如何？
\end{exercise}
% SUPPLEMENT: id=v14-ch06-ex05-solution; parent=ch06-ex05; kind=solution+motivation+hypotheses; origin=independent-derivation
\begin{solution}
\textbf{两个等价关系及无穷端点（重构新增）。}
先补证原提示中的不交集合构造。若存在任意小的正测度集合，给任意正数序列 $\eta_n$，递归选可测 $F_n$，使 $0<\mu(F_1)<\eta_1$，且
\begin{align*}
0<\mu(F_{n+1})<\min\bigl(\eta_{n+1},\mu(F_n)/4\bigr).
\end{align*}
置 $E_n=F_n\setminus\bigcup_{k>n}F_k$。这些集合两两不交，且
\begin{align*}
0<\tfrac23\mu(F_n)\leqslant\mu(E_n)\leqslant\mu(F_n)<\eta_n,
\end{align*}
因为后续测度之和小于 $\mu(F_n)\sum_{j\geqslant1}4^{-j}=\mu(F_n)/3$。若存在任意大的有限测度集合，则在选出有限多个不交有限测度集后，可取有限测度 $F$ 满足 $\mu(F)>\mu(\bigcup_{j<n}E_j)+1$，再令 $E_n=F\setminus\bigcup_{j<n}E_j$，就有 $1<\mu(E_n)<\infty$。

对第一等价关系的“若”方向，令 $\theta=1-p/q>0$，用上述小集合构造取 $m_n=\mu(E_n)<2^{-n/\theta}$，并定义 $a_n=m_n^{-1/q}$。不交性给
\begin{align*}
\int\left|\sum_na_n\chi_{E_n}\right|^p=\sum_nm_n^\theta<\infty,
\qquad
\int\left|\sum_na_n\chi_{E_n}\right|^q=\sum_n1=\infty.
\end{align*}
函数在每个点至多有一个非零项，故是有限复值可测函数。反之，若不存在任意小正测度集合，则有 $\delta>0$，使所有可测正测度集合的测度至少为 $\delta$。对 $f\in L^p$，若 $t>\delta^{-1/p}\|f\|_p$，则
\begin{align*}
\mu(|f|>t)\leqslant t^{-p}\|f\|_p^p<\delta,
\end{align*}
所以该集合必为零集。于是 $\|f\|_\infty\leqslant\delta^{-1/p}\|f\|_p$。若 $f$ 非零，$\int|f|^q\leqslant\|f\|_\infty^{q-p}\int|f|^p<\infty$；零函数也成立，故 $L^p\subseteq L^q$。

对第二等价关系的“若”方向，取上述大集合构造，使 $m_n=\mu(E_n)\geqslant1$，令 $a_n=(nm_n)^{-1/p}$。则
\begin{align*}
\int|f|^p=\sum_n\frac1n=\infty,
\qquad
\int|f|^q=\sum_nn^{-q/p}m_n^{1-q/p}
\leqslant\sum_nn^{-q/p}<\infty.
\end{align*}
反之，若所有有限测度集合的测度都有统一有限上界 $M$，给 $f\in L^q$，集合 $A_n=\{|f|>1/n\}$ 各有有限测度，所以 $\mu(A_n)\leqslant M$。由测度从下连续性，$\mu(\{f\ne0\})\leqslant M$。在这个有限测度集合上应用有限测度包含命题，得 $f\in L^p$；若支集零测度则结论直接成立。

当 $q=\infty$ 时，第一个等价关系仍然正确。“只有若”已由 $\delta$ 下界论证证明；“若”方向取不交 $E_n$ 满足 $0<\mu(E_n)<2^{-n}n^{-p}$，令 $f=\sum_nn\chi_{E_n}$，则 $\int|f|^p<\infty$ 而 $\|f\|_\infty=\infty$。第二个关系应改为
\begin{align*}
L^\infty\not\subseteq L^p\quad\Longleftrightarrow\quad\mu(X)=\infty.
\end{align*}
总测度有限时正文包含定理成立；总测度无穷时常数1属于 $L^\infty$ 却不属于 $L^p$。此时“任意大有限测度集合”不再必要：单点空间上令唯一非空集测度为无穷，则没有正有限测度集合，但仍有上述不包含。
\end{solution}
\begin{remark}
\textbf{动机与条件说明（重构新增）。}小测度集允许极高尖峰，大有限测度集允许很低的值累积成发散质量。不交性使积分精确拆成数列和，不能只挑重叠小集便直接相加。$0<p<q<\infty$ 提供 $\theta>0$ 和 $q/p>1$ 两个可和性条件；无需 Banach 理论，所以允许 $p<1$。有限 $q$ 还迫使非零值域位于有限测度水平集的可数并内；$q=\infty$ 没有这一约束，正是第二判据改变的原因。
\end{remark}
% END SUPPLEMENT: id=v14-ch06-ex05-solution
% SOURCE p=202; printed=187; object=ch06-ex06; number=6; region=upper
\begin{exercise}\label{exercise:ch06-integrability-intervals}
设 $0<p_0<p_1<\infty$。在 $(0,\infty)$ 的 Lebesgue 测度下找函数 $f$，使 $f\in L^p$ 当且仅当：
\begin{enumerate}
\item\label{item:ch06-integrability-open} $p_0<p<p_1$；
\item\label{item:ch06-integrability-closed} $p_0\leqslant p\leqslant p_1$；
\item\label{item:ch06-integrability-singleton} $p=p_0$。
\end{enumerate}
（考虑形如 $f(x)=x^{-a}|\ln x|^b$ 的函数。）
\end{exercise}
% SUPPLEMENT: id=v14-ch06-ex06-solution; parent=ch06-ex06; kind=solution+motivation+hypotheses; origin=independent-derivation
\begin{solution}
\textbf{三类构造与端点验证（重构新增）。}原题规定 $p_1<\infty$；下面同时补充 $p_1=\infty$ 的推广分支，并明确区分它与原题。所有函数在 $[e^{-1},e]$ 上取零；以下只规定左段 $(0,e^{-1})$ 与右段 $(e,\infty)$。先记两个判别。对 $a>0$ 和实数 $b$，代换 $t=-\ln x$ 或 $t=\ln x$，分别得到
\begin{align*}
\int_0^{e^{-1}}x^{-ap}|\ln x|^{bp}\,\mathrm dx
 &=\int_1^\infty e^{(ap-1)t}t^{bp}\,\mathrm dt,\\
\int_e^\infty x^{-ap}(\ln x)^{bp}\,\mathrm dx
 &=\int_1^\infty e^{(1-ap)t}t^{bp}\,\mathrm dt.
\end{align*}
指数系数负时积分有限，正时发散：因为对每个 $\eta>0$，$\ln t/t\to0$ 给出充分大 $t$ 时 $t^{|bp|}\leqslant e^{\eta t}$，从而可与正或负指数比较。系数为零时由幂积分判据，恰在 $bp<-1$ 时有限。
\begin{enumerate}
\item 若 $p_1<\infty$，左段取 $x^{-1/p_1}$，右段取 $x^{-1/p_0}$。左段要求 $p<p_1$，右段要求 $p>p_0$；两段都有非负积分，因此全域可积恰为两个条件的交，且左段无界排除 $p=\infty$。若 $p_1=\infty$，左段改为 $-\ln x$，右段不变。左段对每个有限 $p>0$ 的积分为 $\int_1^\infty e^{-t}t^p\,\mathrm dt<\infty$，但本质无界，故恰保留 $p_0<p<\infty$。
\item 若 $p_1<\infty$，左段取
$x^{-1/p_1}|\ln x|^{-2/p_1}$，右段取
$x^{-1/p_0}(\ln x)^{-2/p_0}$。上述判别给左段恰在 $p\leqslant p_1$ 可积，右段恰在 $p\geqslant p_0$ 可积，临界积分均为 $\int_1^\infty t^{-2}\,\mathrm dt$。左段本质无界，故不额外含 $\infty$。若 $p_1=\infty$，左段改为常数1，右段仍取上述含对数的函数；两段都处处有界，因此也属于 $L^\infty$，有限指数则恰满足 $p\geqslant p_0$。
\item 左段取 $x^{-1/p_0}|\ln x|^{-2/p_0}$，右段取 $x^{-1/p_0}(\ln x)^{-2/p_0}$。前者恰允许有限 $p\leqslant p_0$，后者恰允许有限 $p\geqslant p_0$，故只有 $p=p_0$。左段无界，排除无穷指数。
\end{enumerate}
所有分段函数可测，在每个 $x>0$ 取有限值。以上声称左段无界之处，代换后其大小为正指数增长乘一个固定幂，趋于无穷；连续性保证每个高阈值被超过的集合有正测度，故确为本质无界。
\end{solution}
\begin{remark}
\textbf{动机与条件说明（重构新增）。}原点控制允许指数的上端，无穷远控制下端；对数因子只改变临界指数是否纳入，不改变严格幂次边界。把中间区间置零避免 $\ln1=0$ 造成无关奇性。$p_0>0$ 使构造幂次有意义，$p_0<p_1$ 使预期区间非空；无穷上端须另用有界性或本质无界性核对，不能只计算有限 $p$ 的积分。
\end{remark}
% END SUPPLEMENT: id=v14-ch06-ex06-solution
% SOURCE p=202; printed=187; object=ch06-ex07; number=7; region=upper
\begin{exercise}\label{exercise:ch06-large-exponent-limit}
若对某个 $p<\infty$ 有 $f\in L^p\cap L^\infty$，从而对所有 $q>p$ 都有 $f\in L^q$，则 $\|f\|_\infty=\lim_{q\to\infty}\|f\|_q$。
\end{exercise}
% SUPPLEMENT: id=v14-ch06-ex07-solution; parent=ch06-ex07; kind=solution+motivation+hypotheses; origin=independent-derivation
\begin{solution}
\textbf{上下界夹逼（重构新增）。}这里有限指数的定义域为 $0<p<\infty$。记 $M=\|f\|_\infty$。若 $M=0$，则 $f=0$ 几乎处处，所有范数都为零。若 $M>0$，对 $q>p$ 有
\begin{align*}
\int|f|^q\leqslant M^{q-p}\|f\|_p^p,
\qquad \|f\|_q\leqslant M^{1-p/q}\|f\|_p^{p/q}.
\end{align*}
因此 $\limsup_{q\to\infty}\|f\|_q\leqslant M$。任取 $0<a<M$，本质上确界的定义给 $\mu(E_a)>0$，其中 $E_a=\{|f|>a\}$；又因 $a^p\mu(E_a)\leqslant\|f\|_p^p$，此测度有限。故
\begin{align*}
\|f\|_q\geqslant a\mu(E_a)^{1/q}\longrightarrow a.
\end{align*}
先取下极限，再令 $a\uparrow M$，得到反向下界，因而极限等于 $M$。
\end{solution}
\begin{remark}
\textbf{动机与条件说明（重构新增）。}大指数越来越强调函数接近最大模长的部分；固定正有限测度集的测度因子在开 $q$ 次方后趋于1。某个有限 $p$ 的可积性使这些水平集有限，并控制所有更大指数的积分；本质有界性提供统一上界。本题无需 $p\geqslant1$、有限总测度或本质最大值在某点取得。若删去有限指数可积性，无穷测度空间上的常数1会使全部有限 $q$ 范数为无穷，而无穷范数为1。
\end{remark}
% END SUPPLEMENT: id=v14-ch06-ex07-solution
% SOURCE p=202; printed=187; object=ch06-ex08; number=8; region=upper-middle
\begin{exercise}\label{exercise:ch06-zero-exponent-limit}
设 $\mu(X)=1$，且对某个 $p>0$ 有 $f\in L^p$，从而 $0<q<p$ 时 $f\in L^q$。
\begin{enumerate}
\item\label{item:ch06-zero-exponent-jensen} $\ln\|f\|_q\geqslant\int\ln|f|$。（在\subcref{exercise:ch03-convex-jensen}{item:ch03-jensen}中取 $F(t)=e^t$。）
\item\label{item:ch06-zero-exponent-difference} $(\int|f|^q-1)/q\geqslant\ln\|f\|_q$，且当 $q\to0$ 时，$(\int|f|^q-1)/q\to\int\ln|f|$。
\item\label{item:ch06-zero-exponent-geometric} $\lim_{q\to0}\|f\|_q=\exp(\int\ln|f|)$。
\end{enumerate}
\end{exercise}
% SUPPLEMENT: id=v14-ch06-ex08-solution; parent=ch06-ex08; kind=solution+motivation+hypotheses; origin=independent-derivation
\begin{solution}
\textbf{三个结论及扩展实数情形（重构新增）。}所有极限均从 $q>0$ 一侧取。采用 $\ln0=-\infty$、$e^{-\infty}=0$。先设所给的 $p$ 有限；若已知的指数为无穷，可由 $\mu(X)=1$ 先选任意正有限指数。因为 $t\geqslant1$ 时 $\ln t\leqslant t^p/p$，所以 $L=\int\ln|f|$ 的正部有限，$L$ 有定义且属于 $[-\infty,\infty)$。又有 $|f|^q\leqslant1+|f|^p$，保证以下积分有限。
\begin{enumerate}
\item 若 $L=-\infty$，不等式按扩展实数次序直接成立。若 $L$ 有限，则 $Y=\ln|f|$ 是实值可积函数（在可测零集上任意补有限值）。由指数函数的切线不等式 $e^u\geqslant1+u$，积分得到
\begin{align*}
e^{-qL}\int|f|^q
=\int e^{q(Y-L)}\geqslant1+q\int(Y-L)=1.
\end{align*}
因此 $q^{-1}\ln\int|f|^q\geqslant L$，即所求；这也是原提示的 Jensen 不等式在此情形的直接证明。
\item 令 $A_q=\int|f|^q$。若 $A_q>0$，由 $\ln A_q\leqslant A_q-1$ 得
\begin{align*}
\ln\|f\|_q=\frac{\ln A_q}{q}\leqslant\frac{A_q-1}{q}.
\end{align*}
若 $A_q=0$，左侧为 $-\infty$，不等式仍成立。为证明右侧的极限，利用 $\mu(X)=1$ 把它写成 $\int (|f|^q-1)/q$。在 $t>1$ 上，凸函数 $q\mapsto t^q$ 的割线斜率 $(t^q-1)/q$ 随 $q$ 增大而增大。因此对 $0<q\leqslant p/2$，
\begin{align*}
0\leqslant\frac{t^q-1}{q}\leqslant\frac{2}{p}(t^{p/2}-1),
\end{align*}
且该商趋于 $\ln t$。右端代入 $t=|f|$ 后可积，控制收敛给出 $\{|f|>1\}$ 上的积分极限。对 $0\leqslant t\leqslant1$，$ (1-t^q)/q$ 随 $q\downarrow0$ 增大，趋于 $-\ln t$，包括 $t=0$ 时趋于无穷。单调收敛定理因而给出 $\{|f|\leqslant1\}$ 上负部积分的极限，允许无穷。两部分合并得到
\begin{align*}
\lim_{q\downarrow0}\frac{A_q-1}{q}=L.
\end{align*}
上述单调性先保证沿任意递减到零的序列结论成立，再用斜率的单调性夹逼其余 $q$，故为完整的右极限。
\item 若 $L$ 有限，前两问夹逼 $L\leqslant\ln\|f\|_q\leqslant(A_q-1)/q$，从而对数趋于 $L$。若 $L=-\infty$，第二问的上界趋于 $-\infty$，也给 $\ln\|f\|_q\to-\infty$。取指数分别得到 $e^L$ 或零，包括 $f=0$ 几乎处处的情形。
\end{enumerate}
\end{solution}
\begin{remark}
\textbf{动机与条件说明（重构新增）。}当指数趋零，$t^q=e^{q\ln t}$ 的一阶变化揭示几何平均；不是把 $q=0$ 代入范数公式。总测度1使积分中的常数项恰为1，也是 Jensen 概率平均的归一化。某个正指数的可积性控制对数正部及上方差商，而对数负部不必可积，因此必须保留 $L=-\infty$ 的可能。这里不要求 $f$ 处处为正或远离零；例如 $f$ 在正测度集上为零时，几何平均为零。若总测度不为1，常数函数的 $\|1\|_q=\mu(X)^{1/q}$ 已表明结论不能原样照用。
\end{remark}
% END SUPPLEMENT: id=v14-ch06-ex08-solution
% SOURCE p=202; printed=187; object=ch06-ex09; number=9; region=middle
\begin{exercise}\label{exercise:ch06-norm-measure-convergence}
设 $1\leqslant p<\infty$。若 $\|f_n-f\|_p\to0$，则 $f_n\to f$ 依测度收敛，因此某个子列几乎处处收敛到 $f$。另一方面，若 $f_n\to f$ 依测度收敛，且对所有 $n$ 都有 $|f_n|\leqslant g\in L^p$，则 $\|f_n-f\|_p\to0$。
\end{exercise}
% SUPPLEMENT: id=v14-ch06-ex09-solution; parent=ch06-ex09; kind=solution+motivation+hypotheses; origin=independent-derivation
\begin{solution}
\textbf{两方向的完整论证（重构新增）。}对任意 $\varepsilon>0$，Chebyshev 估计给
\begin{align*}
\mu(|f_n-f|>\varepsilon)\leqslant\varepsilon^{-p}\|f_n-f\|_p^p\longrightarrow0.
\end{align*}
依测度收敛到可测 $f$ 时，可选严格递增的 $n_k$ 使 $\mu(|f_{n_k}-f|>2^{-k})<2^{-k}$。这些坏集的上极限测度为零，因为其每条尾部并的测度至多为 $\sum_{j\geqslant k}2^{-j}\to0$。在上极限以外，最终有 $|f_{n_k}-f|\leqslant2^{-k}$，所以几乎处处收敛。

现在假设共同控制条件成立。选上述子列，在它的几乎处处收敛集和全部控制不等式成立的共同满测集上取极限，可得 $|f|\leqslant g$；于是 $f\in L^p$。若范数差不趋零，存在 $\eta>0$ 和子列 $f_{m_j}$，使 $\|f_{m_j}-f\|_p\geqslant\eta$。这个子列仍依测度收敛到 $f$，因而有进一步子列几乎处处收敛。沿进一步子列，$|f_{m_j}-f|^p\to0$ 且 $|f_{m_j}-f|^p\leqslant2^pg^p\in L^1$，控制收敛定理给范数差趋零，矛盾。因此原序列范数收敛。
\end{solution}
\begin{remark}
\textbf{动机与条件说明（重构新增）。}范数直接控制每个固定高度的坏集；依测度收敛本身则不能排除越来越高、越来越窄的尖峰。共同的 $L^p$ 控制函数排除这种质量集中，并通过“任意子列再抽子列”的反证把几乎处处控制收敛提升为全序列结论。有限指数保证 $g^p$ 的积分控制有效；证明事实上对全部 $0<p<\infty$ 也成立。无需总测度有限。共同控制若仅几乎处处成立，可先取可数个可测例外零集之并。
\end{remark}
% END SUPPLEMENT: id=v14-ch06-ex09-solution
% SOURCE p=202; printed=187; object=ch06-ex10; number=10; region=middle
\begin{exercise}\label{exercise:ch06-ae-norm-convergence}
设 $1\leqslant p<\infty$。若 $f_n,f\in L^p$ 且 $f_n\to f$ 几乎处处，则 $\|f_n-f\|_p\to0$ 当且仅当 $\|f_n\|_p\to\|f\|_p$。（使用\cref{exercise:ch02-generalized-dominated}。）
\end{exercise}
% SUPPLEMENT: id=v14-ch06-ex10-solution; parent=ch06-ex10; kind=solution+motivation+hypotheses; origin=independent-derivation
\begin{solution}
\textbf{广义控制收敛的直接实现（重构新增）。}必要性由反三角不等式
$|\|f_n\|_p-\|f\|_p|\leqslant\|f_n-f\|_p$ 立即得到。反之，设范数收敛，记
\begin{align*}
h_n=|f_n-f|^p,
\qquad k_n=2^{p-1}(|f_n|^p+|f|^p).
\end{align*}
由 $|z+w|^p\leqslant2^{p-1}(|z|^p+|w|^p)$，有 $0\leqslant h_n\leqslant k_n$。几乎处处有 $h_n\to0$、$k_n\to k=2^p|f|^p$；范数收敛还给 $\int k_n\to\int k<\infty$。对非负函数 $k_n-h_n$ 用 Fatou 引理，得到
\begin{align*}
\int k\leqslant\liminf_n\int(k_n-h_n)
=\int k-\limsup_n\int h_n.
\end{align*}
这里减法合法，因为 $\int k_n$ 收敛到有限数，且 $0\leqslant\int h_n\leqslant\int k_n$。故 $\limsup_n\int h_n\leqslant0$，即 $\|f_n-f\|_p\to0$。所用幂次估计由 $t\mapsto t^p$ 的凸性施于两个非负模长得到，$p=1$ 也成立。
\end{solution}
\begin{remark}
\textbf{动机与条件说明（重构新增）。}几乎处处收敛给形状的逐点极限，范数收敛则保证总的 $p$ 次质量没有遗失；两者结合排除迁移或集中的剩余质量。各函数属于 $L^p$，使 Fatou 后的有限数相减合法。$p\geqslant1$ 使必要性可用反三角不等式，$p<\infty$ 使范数对应实际积分。没有几乎处处收敛，仅有范数相等并不足够，例如非零 $f$ 与恒为 $-f$ 的序列。无穷指数也不能原样推广：在 $(0,1)$ 上取 $f=1$、$f_n=1-\chi_{(0,1/n)}$，则几乎处处收敛且无穷范数均为1（$n\geqslant2$），但差的无穷范数恒为1。
\end{remark}
% END SUPPLEMENT: id=v14-ch06-ex10-solution
% SOURCE p=202; printed=187; object=ch06-ex11; number=11; region=middle-lower
\begin{exercise}\label{exercise:ch06-essential-range}
设 $f$ 为 $X$ 上的可测函数。
\begin{definition*}[本质值域]\phantomsection\label{definition:ch06-essential-range-local}
$f$ 的\textbf{本质值域} $R_f$ 是所有满足下述条件的 $z\in\mathbb C$ 组成的集合：对每个 $\varepsilon>0$，集合 $\{x:|f(x)-z|<\varepsilon\}$ 都有正测度。
\end{definition*}
% SUPPLEMENT: id=v14-ch06-essential-range-local-remark; parent=ch06-ex11; kind=motivation+hypotheses; origin=independent-derivation
\begin{remark}
\textbf{局部定义的动机与条件（重构新增）。}本质值域保留在任意精度下都能在正测度集合上观察到的数值，舍弃只出现在零集上的值。可测性保证圆盘原像可测；“每个正半径”不能改成只查一个邻域，否则会把远离实际值域的点纳入。定义只要求邻域原像正测度，不要求 $f$ 真正取到该值，也不要求该测度有限。定义对几乎处处相等的代表不变，因为对应原像只可能在它们不相等的可测零集内不同。零测度空间的本质值域为空，须在涉及最大值时另行处理。
\end{remark}
% END SUPPLEMENT: id=v14-ch06-essential-range-local-remark
\begin{enumerate}
\item\label{item:ch06-essential-range-closed} $R_f$ 是闭集。
\item\label{item:ch06-essential-range-compact} 若 $f\in L^\infty$，则 $R_f$ 紧，且 $\|f\|_\infty=\max\{|z|:z\in R_f\}$。
\end{enumerate}
\end{exercise}
% SUPPLEMENT: id=v14-ch06-ex11-solution; parent=ch06-ex11; kind=solution+motivation+hypotheses+source-qualification; origin=independent-derivation
\begin{remark}
\textbf{原题的退化情形说明（重构新增）。}原题第二问的最大值等式需要 $\mu(X)>0$。若 $\mu(X)=0$，则 $R_f=\varnothing$ 仍是紧集，$\|f\|_\infty=0$，但空集上的最大值不存在。下面保留原题并在非零测度情形证明该等式；同时适用于两种情形的写法为 $\|f\|_\infty=\sup(\{0\}\cup\{|z|:z\in R_f\})$。
\end{remark}
\begin{solution}
\textbf{闭性、紧性与范数（重构新增）。}
\begin{enumerate}
\item 若 $z\notin R_f$，存在 $\varepsilon>0$ 使 $\mu(f^{-1}(B(z,\varepsilon)))=0$。当 $|w-z|<\varepsilon/2$ 时，$B(w,\varepsilon/2)\subseteq B(z,\varepsilon)$，故 $w\notin R_f$。所以补集开，$R_f$ 闭。
\item 先证明 $f(x)\in R_f$ 几乎处处。复平面有由有理实虚部中心和正有理半径构成的可数基。对每个 $z\notin R_f$，取原像零测的圆盘邻域，再取包含 $z$ 且落在该圆盘中的基元 $U$。所有具有零测原像的这类基元的并恰为 $\mathbb C\setminus R_f$：反向是因为一个点若属于此开基元，就有更小的零测原像圆盘。于是 $f^{-1}(\mathbb C\setminus R_f)$ 是可数个可测零集的并，为零集。

设 $M=\|f\|_\infty<\infty$。若 $|z|>M$，取半径 $0<\varepsilon<|z|-M$，则 $f(x)\in B(z,\varepsilon)$ 迫使 $|f(x)|>M$，而此集合零测。因此 $R_f\subseteq\{|z|\leqslant M\}$，结合第一问，它是闭有界的复平面子集，故紧。当 $\mu(X)>0$ 时，上段的几乎处处包含关系还保证 $R_f\ne\varnothing$。连续函数 $z\mapsto|z|$ 在其上取得最大值 $A\leqslant M$；又因 $|f|\leqslant A$ 几乎处处，$M\leqslant A$，得到等式。若 $\mu(X)=0$，则定义直接给 $R_f=\varnothing$ 和补充说明中的统一上确界公式。
\end{enumerate}
\end{solution}
\begin{remark}
\textbf{结论的动机与条件说明（重构新增）。}本质上确界可以不由函数在某个点实际取得，却会在非空紧本质值域上成为真正最大模长。本质有界性提供有界，第一问提供闭，复平面的有限维性质把两者合成紧性。证明 $f$ 几乎处处落在本质值域时关键使用值域的可数基，不能直接把任意多个零集取并。无需测度空间有限、半有限或完备；$\mu(X)>0$ 的作用只是排除空本质值域导致的最大值无定义。
\end{remark}
% END SUPPLEMENT: id=v14-ch06-ex11-solution
% SOURCE p=202; printed=187; object=ch06-ex12; number=12; region=lower
\begin{exercise}\label{exercise:ch06-inner-product-exponent}
若 $p\ne2$，则 $L^p$ 范数不是由 $L^p$ 上的某个内积导出的，除非处于 $\dim(L^p)\leqslant1$ 的平凡情形。（证明平行四边形恒等式不成立。）
\end{exercise}
% SUPPLEMENT: id=v14-ch06-ex12-solution; parent=ch06-ex12; kind=solution+motivation+hypotheses; origin=independent-derivation
\begin{solution}
\textbf{先找到两个不交方向（重构新增）。}先按本节赋范空间的范围，设 $1\leqslant p\leqslant\infty$。把可测集 $E$ 称为此段中的可用集合，当且仅当 $\chi_E\in L^p$；有限 $p$ 时等价于 $\mu(E)<\infty$，$p=\infty$ 时所有可测集均可用。简单函数稠密性说明这些示性函数的线性包稠密于 $L^p$。

若不存在两个不交、正测度的可用集合，任意两个正测度可用集合 $E,F$ 必须几乎处处相等。否则，$E\setminus F$ 或 $F\setminus E$ 中至少一个有正测度；它仍可用，并分别与 $F$ 或 $E$ 不交，矛盾。因此所有非零可用示性函数都相等为同一个 $\chi_A$（若一个也没有，则稠密性给 $L^p=\{0\}$）。其线性包至多一维；有限维子空间在赋范空间中闭，故由稠密性得到 $\dim(L^p)\leqslant1$。取逆否命题，若维数至少为2，就有两个不交正测度可用集合 $E,F$。

有限 $p$ 时令 $u=\mu(E)^{-1/p}\chi_E$、$v=\mu(F)^{-1/p}\chi_F$，则
\begin{align*}
\|u\|_p=\|v\|_p=1,
\qquad \|u+v\|_p=\|u-v\|_p=2^{1/p}.
\end{align*}
如果范数由内积导出，平行四边形恒等式迫使 $2\cdot2^{2/p}=4$，即 $p=2$。当 $p=\infty$ 时，直接取 $u=\chi_E$、$v=\chi_F$，四个有关范数都为1，平行四边形恒等式却要求 $2=4$，同样矛盾。

零空间上的唯一内积给零范数；一维情形选非零向量 $e$，定义 $\langle ae,be\rangle=a\overline b\|e\|_p^2$，确实导出原范数，因此所列例外真实存在。
\end{solution}
\begin{remark}
\textbf{动机与条件说明（重构新增）。}内积把正交方向的平方长度相加，而不交支集使 $L^p$ 范数把 $p$ 次方相加；两种规则只有 $p=2$ 相合。维数条件保证能找到两个独立的不交方向，不能在任意测度空间上直接假定两个正有限测度集存在。有限 $p$ 与无穷 $p$ 的可用集合不同，所以无穷端点须另证。无需半有限性或总测度有限。对 $0<p<1$，原量一般已不是范数；若存在两个不交正有限测度集，正文的三角不等式反例也立即排除它由内积导出。
\end{remark}
% END SUPPLEMENT: id=v14-ch06-ex12-solution
% SOURCE p=202; printed=187; object=ch06-ex13; number=13; region=lower
\begin{exercise}\label{exercise:ch06-lebesgue-separability}
当 $1\leqslant p<\infty$ 时，$L^p(\mathbb R^n,m)$ 可分。但是 $L^\infty(\mathbb R^n,m)$ 不可分。（存在不可数集 $\mathscr F\subseteq L^\infty$，使不同的 $f,g\in\mathscr F$ 都满足 $\|f-g\|_\infty\geqslant1$。）
\end{exercise}
% SUPPLEMENT: id=v14-ch06-ex13-solution; parent=ch06-ex13; kind=solution+motivation+hypotheses; origin=independent-derivation
\begin{solution}
\textbf{可数稠密族与不可数分离族（重构新增）。}取 $n\geqslant1$。设 $\mathcal A$ 为所有有限个有界有理端点开长方体的并，以及空集所组成的族。长方体族可数，有限列的集合仍可数，故 $\mathcal A$ 可数。设 $\mathcal D$ 是所有有限和 $\sum_{j=1}^Na_j\chi_{A_j}$，其中 $a_j\in\mathbb Q+i\mathbb Q$、$A_j\in\mathcal A$；它同样可数。

给任意有限测度的 Lebesgue 可测集 $E$ 及 $\delta>0$，正则性提供紧集 $K\subseteq E$ 和开集 $O\supseteq E$，满足 $m(E\setminus K)<\delta/2$、$m(O\setminus E)<\delta/2$。每个 $x\in K$ 有包含它且落在 $O$ 内的有界有理端点开长方体；紧性选出有限子覆盖，它们的并 $A\in\mathcal A$ 满足 $K\subseteq A\subseteq O$，于是 $m(E\mathbin{\triangle}A)<\delta$。若 $K$ 为空，可取 $A=\varnothing$，同样成立。因此
\begin{align*}
\|\chi_E-\chi_A\|_p=m(E\mathbin{\triangle}A)^{1/p}
\end{align*}
可任意小。有限测度支集的简单函数在 $L^p$ 中稠密，逐项用上述 $A$ 逼近其示性函数，再把复系数逼近为有理实虚部数，即得 $\mathcal D$ 稠密。具体地，每项系数误差的范数为 $|c-a|m(A)^{1/p}$，而有限项之和由三角不等式控制，所以两层逼近都能达到任意指定误差。

对无穷指数，令 $f_t=\chi_{(0,t)\times(0,1)^{n-1}}$，$0<t<1$；$n=1$ 时取区间 $(0,t)$。不同参数对应的差在正测度长方体上模长为1，因此 $\|f_s-f_t\|_\infty=1$。这是不可数的分离族。若存在可数稠密集 $D$，则以各 $f_t$ 为中心、半径 $1/3$ 的两两不交开球都须包含 $D$ 的点；把每个球对应到固定枚举中落入它的第一个点，得到不可数族到可数集的单射，矛盾。
\end{solution}
\begin{remark}
\textbf{动机与条件说明（重构新增）。}有限指数允许忽略小测度误差，因而可用有理几何数据编码稠密族；无穷范数对任何正测度误差仍收取完整高度。Lebesgue 测度的内外正则性和欧氏空间的可数基是此构造的依据，结论不是任意测度空间都成立。原题默认正整数维数：若 $n=0$，空间只有一个点，$L^\infty$ 一维，不可分结论不适用。有限指数允许先用有限测度支集的简单函数逼近，无需 $\mathbb R^n$ 的总测度有限。
\end{remark}
% END SUPPLEMENT: id=v14-ch06-ex13-solution
% SOURCE p=202; printed=187; object=ch06-ex14; number=14; region=lower
\begin{exercise}\label{exercise:ch06-multiplication-operator}
若 $g\in L^\infty$，则由 $Tf=fg$ 定义的算子在 $L^p$ 上有界，其中 $1\leqslant p\leqslant\infty$。其算子范数至多为 $\|g\|_\infty$；若 $\mu$ 半有限，则等号成立。
\end{exercise}
% SUPPLEMENT: id=v14-ch06-ex14-solution; parent=ch06-ex14; kind=solution+motivation+hypotheses; origin=independent-derivation
\begin{solution}
\textbf{上下界及半有限性的用途（重构新增）。}令 $M=\|g\|_\infty$。乘积可测，改变任一代表只会在可测零集上改变乘积；线性直接来自逐点运算。$|fg|\leqslant M|f|$ 几乎处处，有限指数取积分，无穷指数取本质上确界，都给 $\|Tf\|_p\leqslant M\|f\|_p$。故 $T$ 良定义且有界，$\|T\|\leqslant M$。

若 $M=0$，算子为零，已得等号。设 $M>0$ 且 $p<\infty$。对 $0<\varepsilon<M$，集合 $A=\{|g|>M-\varepsilon\}$ 有正测度。半有限性保证存在可测 $E\subseteq A$，满足 $0<\mu(E)<\infty$：若 $A$ 自身测度有限可直接取它，若无穷则用半有限性的定义。取 $f=\mu(E)^{-1/p}\chi_E$，则 $\|f\|_p=1$ 且
\begin{align*}
\|Tf\|_p^p=\frac1{\mu(E)}\int_E|g|^p\geqslant(M-\varepsilon)^p.
\end{align*}
所以 $\|T\|\geqslant M-\varepsilon$；令 $\varepsilon\downarrow0$ 得等号。若 $p=\infty$，在 $M>0$ 时必有 $\mu(X)>0$，常数1具有范数1，而且 $T1=g$，直接给 $\|T\|\geqslant M$，此端点无需半有限性。
\end{solution}
\begin{remark}
\textbf{动机与条件说明（重构新增）。}乘法算子的最大放大率应当由乘子接近本质最大值的区域检测出来。对有限指数，测试函数必须属于 $L^p$，所以要从该区域选出正有限测度部分；半有限性正提供这一步。它是所给一般等号结论的充分条件，不是每个具体乘子都必需的条件。若测度空间只有一个点且其测度为无穷，则所有有限 $p$ 的 $L^p$ 都是零空间；取 $g=1$，有 $\|T\|=0$ 而 $\|g\|_\infty=1$，表明不能无条件删去这项假设。算子范数在零空间上以单位闭球的上确界定义为零。无穷端点由常数测试解决，不受此障碍影响。
\end{remark}
% END SUPPLEMENT: id=v14-ch06-ex14-solution
% SOURCE p=202; printed=187; object=ch06-ex15; number=15; region=bottom
\begin{exercise}[Vitali 收敛定理]\label{exercise:ch06-vitali-convergence}
设 $1\leqslant p<\infty$，且 $\{f_n\}_{n=1}^\infty\subseteq L^p$。$\{f_n\}$ 关于 $L^p$ 范数为 Cauchy 列的充要条件是以下三条成立：
\begin{enumerate}
\item\label{item:ch06-vitali-measure} $\{f_n\}$ 依测度为 Cauchy 列；
\item\label{item:ch06-vitali-integrability} $\{|f_n|^p\}$ 一致可积，见\cref{exercise:ch03-uniform-integrability}；
\item\label{item:ch06-vitali-tail} 对每个 $\varepsilon>0$，存在 $E\subseteq X$，使 $\mu(E)<\infty$ 且对所有 $n$ 都有 $\int_{E^c}|f_n|^p<\varepsilon$。
\end{enumerate}
（为证充分性：给定 $\varepsilon>0$，按第三条选 $E$，并令 $A_{mn}=\{x\in E:|f_m(x)-f_n(x)|>\varepsilon\}$。当 $m,n$ 足够大时，$|f_m-f_n|^p$ 在 $E\setminus A_{mn}$、$A_{mn}$ 和 $E^c$ 上的积分都很小，分别基于三种不同原因。）
\end{exercise}
% SUPPLEMENT: id=v14-ch06-ex15-solution; parent=ch06-ex15; kind=solution+motivation+hypotheses; origin=independent-derivation
\begin{solution}
\textbf{三个条件的必要性和充分性（重构新增）。}本题的一致可积性沿所引习题的约定，指对每个 $\eta>0$ 存在 $\delta>0$，使可测 $A$ 满足 $\mu(A)<\delta$ 时，所有 $n$ 都有 $\int_A|f_n|^p<\eta$。被积函数非负，因此该习题中积分外的绝对值在此没有区别。记 $C=2^{p-1}$，有 $|a-b|^p\leqslant C(|a|^p+|b|^p)$。

先证必要性。由 $L^p$ 完备性，范数 Cauchy 列收敛到某个 $f\in L^p$。第一条可直接由
\begin{align*}
\mu(|f_m-f_n|>t)\leqslant t^{-p}\|f_m-f_n\|_p^p
\end{align*}
得到。为证第二条，先说明任意 $u\in L^1$ 的积分对小测度集合绝对连续。给 $\eta>0$，选择 $K>0$ 使 $\int_{\{|u|>K\}}|u|<\eta/2$，再取 $\delta<\eta/(2K)$，则 $\mu(A)<\delta$ 蕴含 $\int_A|u|\leqslant K\mu(A)+\int_{\{|u|>K\}}|u|<\eta$。尾积分可以趋零，是控制收敛定理的直接结果。

给定本题的 $\eta>0$，取 $N$ 使 $n\geqslant N$ 时 $\|f_n-f\|_p^p<\eta/(2C)$。按上段为 $|f|^p$ 选小测度阈值，使 $C\int_A|f|^p<\eta/2$；再缩小该阈值，使有限多个 $|f_1|^p,\ldots,|f_{N-1}|^p$ 的积分都小于 $\eta$。于是对 $n\geqslant N$，
\begin{align*}
\int_A|f_n|^p\leqslant C\int_A|f|^p+C\|f_n-f\|_p^p<\eta,
\end{align*}
有限个先前项也满足要求，第二条成立。

为证第三条，先对任意 $h\in L^p$ 取 $E_k(h)=\{|h|>1/k\}$。其测度不超过 $k^p\|h\|_p^p$，因而有限；同时 $\int_{E_k(h)^c}|h|^p\to0$，因为被积函数逐点趋零并由 $|h|^p$ 控制。给 $\varepsilon>0$，取 $N$ 使 $n\geqslant N$ 时 $\|f_n-f\|_p^p<\varepsilon/(2C)$。为 $f$ 选择有限测度集 $E_0$ 使外部积分小于 $\varepsilon/(2C)$，为每个 $j<N$ 选择有限测度集 $E_j$ 使 $\int_{E_j^c}|f_j|^p<\varepsilon$。它们的有限并 $E$ 仍有有限测度；对 $n\geqslant N$ 用同一幂次估计得到 $\int_{E^c}|f_n|^p<\varepsilon$，前面有限项也成立。

再证充分性。给目标误差 $\varepsilon>0$，令 $\eta=\varepsilon^p/(2(1+4C))$。按第三条取有限测度 $E$，使所有项在 $E^c$ 上的 $p$ 次积分小于 $\eta$；按第二条取 $\delta>0$，使测度小于 $\delta$ 的集合上这些积分都小于 $\eta$。选 $t>0$ 满足 $t^p\mu(E)<\eta$；若 $\mu(E)=0$，任意正 $t$ 即可。由第一条，当 $m,n$ 足够大时，集合 $A=E\cap\{|f_m-f_n|>t\}$ 满足 $\mu(A)<\delta$。分三部分积分，得到
\begin{align*}
\int_{E\setminus A}|f_m-f_n|^p&\leqslant t^p\mu(E)<\eta,\\
\int_A|f_m-f_n|^p&\leqslant C\left(\int_A|f_m|^p+\int_A|f_n|^p\right)<2C\eta,\\
\int_{E^c}|f_m-f_n|^p&\leqslant C\left(\int_{E^c}|f_m|^p+\int_{E^c}|f_n|^p\right)<2C\eta.
\end{align*}
所以 $\|f_m-f_n\|_p^p<(1+4C)\eta<\varepsilon^p$，即范数 Cauchy 性。
\end{solution}
\begin{remark}
\textbf{动机与逐条件说明（重构新增）。}第一条控制固定阈值的坏集有多大，第二条排除质量集中在越来越小的集合上，第三条排除质量逃向任何固定有限测度区域之外。三者各有独立用途。在 $(0,1)$ 上，$f_n=n^{1/p}\chi_{(0,1/n)}$ 依测度趋零且第三条自动成立，但第二条失败，$L^p$ 范数恒为1而不能 Cauchy。在 $\mathbb R$ 上，$f_n=\chi_{(n,n+1)}/n^{1/p}$ 范数趋零，说明移动支集本身不一定造成障碍；真正显示第三条不可删的例子是 $f_n=n^{-1/p}\chi_{(0,n)}$：它依测度趋零，因 $|f_n|^p\leqslant1$ 而一致可积，但对任意有限测度 $E$，$\int_{E^c}|f_n|^p\geqslant1-\mu(E)/n\to1$。第一条也不可删，例如在概率空间上交替取常数0与1，后两条成立却不是依测度 Cauchy。有限 $p$ 把范数化为积分；$p\geqslant1$ 提供已证的 Banach 完备性。无需测度空间本身 $\sigma$ 有限，第三条只对这一个函数族提出共同的有限测度截断。
\end{remark}
% END SUPPLEMENT: id=v14-ch06-ex15-solution
% SOURCE p=202,203; printed=187,188; object=ch06-ex16; number=16; region=bottom-top
\begin{exercise}\label{exercise:ch06-small-p-completeness}
若 $0<p<1$，则公式 $\rho(f,g)=\int|f-g|^p$ 在 $L^p$ 上定义一个度量，使 $L^p$ 成为完备的拓扑向量空间。（在\cref{theorem:ch06-lp-completeness}的证明中，把 $\|f\|_p$ 换为 $\int|f|^p$，其论证对 $p<1$ 仍然成立，因为它只使用三角不等式，而不用范数的齐次性。）
\end{exercise}
% SUPPLEMENT: id=v14-ch06-ex16-solution; parent=ch06-ex16; kind=solution+motivation+hypotheses; origin=independent-derivation
\begin{solution}
\textbf{度量、连续运算及完备性（重构新增）。}若 $a,b\geqslant0$，则
\begin{align*}
(a+b)^p\leqslant a^p+b^p.
\end{align*}
当 $a+b>0$ 时，除以 $(a+b)^p$ 后只需用 $s^p\geqslant s$ 于 $s=a/(a+b)$ 和 $1-s$；当 $a=b=0$ 时等式直接成立。结合复数模的三角不等式，得到 $|z+w|^p\leqslant|z|^p+|w|^p$。

因此 $\rho$ 在 $L^p$ 上有限，对称，且满足三角不等式。其零值条件等价于几乎处处相等，所以在已识别代表的 $L^p$ 上为真正度量。它有平移不变性以及 $\rho(af,ag)=|a|^p\rho(f,g)$。进一步，
\begin{align*}
\rho(f+h,g+k)&\leqslant\rho(f,g)+\rho(h,k),\\
\rho(af,bg)&\leqslant|a|^p\rho(f,g)+|a-b|^p\rho(g,0).
\end{align*}
第一式证明加法联合连续；若 $a_n\to a$ 且 $f_n\to f$ 关于 $\rho$ 收敛，在第二式中取 $(a,f,b,g)=(a_n,f_n,a,f)$，利用 $|a_n|$ 有界便得 $a_nf_n\to af$。定义域与值域均为度量空间，序列连续性等价于连续性，故标量乘法也联合连续。于是为拓扑向量空间。

设 $\{f_n\}$ 为 $\rho$-Cauchy 列。取子列 $\{f_{n_k}\}$，满足 $\rho(f_{n_{k+1}},f_{n_k})<2^{-k}$，并选可测有限值代表。记 $d_k=f_{n_{k+1}}-f_{n_k}$。由单调收敛定理，
\begin{align*}
\int\sum_{k=1}^\infty|d_k|^p=\sum_{k=1}^\infty\rho(f_{n_{k+1}},f_{n_k})<\infty.
\end{align*}
所以可测集合 $N=\{\sum_k|d_k|^p=\infty\}$ 为零集。对 $x\notin N$，最终有 $|d_k(x)|\leqslant1$，从而 $|d_k(x)|\leqslant|d_k(x)|^p$；因此 $\sum_k|d_k(x)|<\infty$。在 $N^c$ 上定义 $f=f_{n_1}+\sum_kd_k$，在 $N$ 上令 $f=0$。这是有限复值可测函数。幂次三角不等式及部分和极限给
\begin{align*}
|f-f_{n_j}|^p\leqslant\sum_{k=j}^\infty|d_k|^p\quad\text{几乎处处},
\qquad
\int|f-f_{n_j}|^p\leqslant\sum_{k=j}^\infty2^{-k}\longrightarrow0.
\end{align*}
取 $j=1$ 并用 $|f|^p\leqslant|f-f_{n_1}|^p+|f_{n_1}|^p$，可知 $f\in L^p$。子列关于 $\rho$ 收敛到 $f$，原列的 Cauchy 性与三角不等式随即给整列收敛：给 $\varepsilon>0$，先使充分后面的任意两项距离小于 $\varepsilon/2$，再在其中选一项子列使它到 $f$ 的距离小于 $\varepsilon/2$。故此度量完备。
\end{solution}
\begin{remark}
\textbf{动机与条件说明（重构新增）。}$p<1$ 破坏了取 $p$ 次方根后的三角不等式，却加强了未开根号的幂次次可加性；因此改用 $\rho$ 仍能建立完整的收敛理论。$p>0$ 保证零值识别和 $|a_n-a|^p\to0$；$p<1$ 用于幂次不等式及“$p$ 次幂可和推出绝对可和”。等价类消除零距离的不同代表，可测零集上的整体补零使构造不依赖测度完备性。度量只有 $p$ 次齐次性，不能把它直接称为范数，也不能未经说明就套用 Banach 空间的范数级数判据；上面的子列级数论证直接证明了所需完备性。
\end{remark}
% END SUPPLEMENT: id=v14-ch06-ex16-solution
\end{document}
